\documentclass[11pt,a4paper]{article}
\usepackage[margin=23mm,headheight=15pt]{geometry}
\usepackage{fontspec}
\setmainfont{P052}[BoldFont={P052-Bold},ItalicFont={P052-Italic},BoldItalicFont={P052-BoldItalic}]
\setsansfont{DejaVu Sans}[Scale=.92]
\setmonofont{DejaVu Sans Mono}[Scale=.80]
\newfontfamily\codemathfont{Latin Modern Math}
\newcommand{\CodeSymbol}[1]{\ifnum"#1="29F8\ensuremath{\big/}\else{\codemathfont\symbol{"#1}}\fi}
\usepackage{amsmath}
\usepackage{unicode-math}
\setmathfont{Latin Modern Math}
\usepackage{xcolor,booktabs,longtable,array,enumitem,microtype}
\PassOptionsToPackage{obeyspaces}{url}
\usepackage{listings,fvextra,xurl,hyperref,bookmark,fancyhdr,titlesec,needspace}
\definecolor{teal}{HTML}{164C51}
\definecolor{muted}{HTML}{516466}
\definecolor{pale}{HTML}{F2F6F6}
\hypersetup{colorlinks=true,linkcolor=teal,urlcolor=teal,pdftitle={Paper II: Lean Guide},pdfauthor={Komorebi P.},bookmarksnumbered=true}
\urlstyle{tt}
\newcommand{\code}[1]{\nolinkurl{#1}}
\newcommand{\source}[1]{\par\smallskip{\footnotesize\color{muted}Source: \code{#1}.}\par\smallskip}
\newcommand{\excerpt}[2]{\Needspace{5\baselineskip}\source{KomorebiVerification/#1.lean}\VerbatimInput[fontsize=\small,frame=leftline,rulecolor=\color{teal},framesep=6pt,breaklines=true,breakanywhere=true,commandchars=\\\{\},xleftmargin=8pt]{excerpts/#2.render.txt}}
\titleformat{\section}{\Large\sffamily\bfseries\color{teal}}{\thesection}{.65em}{}
\titleformat{\subsection}{\large\sffamily\bfseries\color{teal}}{\thesubsection}{.65em}{}
\titleformat{\subsubsection}{\normalsize\sffamily\bfseries\color{teal}}{\thesubsubsection}{.65em}{}
\titlespacing*{\section}{0pt}{2.2ex plus .6ex}{1ex}
\titlespacing*{\subsection}{0pt}{2ex plus .4ex}{.7ex}
\setlength{\parindent}{0pt}
\setlength{\parskip}{.55em}
\setlength{\emergencystretch}{2em}
\setlist{nosep,topsep=.4em,leftmargin=1.5em}
\pagestyle{fancy}\fancyhf{}
\fancyhead[L]{\small\sffamily\color{muted}KOMOREBI \quad PAPER II}
\fancyhead[R]{\small\sffamily\color{muted}LEAN GUIDE}
\fancyfoot[C]{\small\thepage}
\renewcommand{\headrulewidth}{.3pt}
\lstset{basicstyle=\ttfamily\small,backgroundcolor=\color{pale},frame=single,rulecolor=\color{pale},framerule=0pt,xleftmargin=5pt,xrightmargin=5pt,aboveskip=7pt,belowskip=7pt,breaklines=true,breakatwhitespace=false,columns=fullflexible,keepspaces=true,showstringspaces=false,tabsize=2,escapeinside={(*@}{@*)}}
\makeatletter
\renewcommand*\l@subsection{\@dottedtocline{2}{1.5em}{3.5em}}
\makeatother

% Operation Komorebi: shared publication profile.
\usepackage{komorebi-publication}
\renewcommand{\KomorebiTitle}{Projector-Resolved Seed Calculus II}
\renewcommand{\KomorebiSubtitle}{Refinement towers, weighted transfer, and integral separation of operator packets}
\renewcommand{\KomorebiSeries}{Core series II}
\renewcommand{\KomorebiKind}{Lean and evidence guide}
\renewcommand{\KomorebiEdition}{Standardized publication edition / 8 October 2026}
\renewcommand{\KomorebiShortTitle}{Paper II}
\renewcommand{\KomorebiTitleSize}{31}
\renewcommand{\KomorebiTitleLeading}{36}
\renewcommand{\KomorebiProfile}{classical}
\renewcommand{\KomorebiOrientation}{portrait}
\KomorebiSetup
\begin{document}
\KomorebiCover
\KomorebiOpeningNote{A practical reading guide to the complete 25-row contract of the preserved Paper II formal component. The guide explains the source architecture, follows representative proofs, records every mapped target and its boundary, and supplies a reproducible route to rebuilding the formal project.
\medskip

\textbf{Evidence at this edition.} All 22 preserved formal-component files match the release checksum manifest and the baseline copy. The source contains 15 Lean files. The report of 3 September 2026 records 160 mapped declarations and 1,034 audited namespace entries. These are inherited report claims; no raw compiler transcript accompanies this component, and no new Lean compiler run is claimed here.}


\tableofcontents
\clearpage
\section{What this guide verifies and teaches}\label{sec:orientation}

The traditional paper develops the mathematical theory and its written proofs. The visual guide provides finite models for its identities, precision statements, and boundaries. This guide explains how to read and use the preserved Lean component: the types accepted by each theorem, the objects actually constructed, the obligations supplied by a caller, and the evidence attached to the source. These three companions have complementary roles.

The formal component is version 1.1.1, preserved within canonical manuscript release 1.2.1. In the delivered package it is located at \code{source/canonical_v1.2.1/formalization/lean/}. The final traditional edition incorporates a broader written development. Its ordinary proofs of lattice gluing, necessary precision delay, inverse-limit convergence, general DVR arithmetic, holomorphic and exponential defects, and multilevel identities are not converted into Lean proofs by appearing beside formally mapped results.

The fixed formal denominator is 25 rows: 21 numbered canonical statements, one Smith-phase identity, and three mathematical algorithm specifications. The appendix crosswalk identifies their locations in the integrated traditional edition. Throughout this guide, ``row'' means a row of the canonical formal map, never a count of every theorem in the final traditional paper.

\begin{tabular}{p{.24\textwidth}p{.67\textwidth}}
\toprule Quantity & What is counted \\\midrule
15 Lean source files & Four shared foundation modules, eight substantive Paper II modules, two aggregate imports, and one audit module. \\
5,603 source lines & Physical lines in those files, including comments and blank lines. \\
25 map rows & Statement families, constructions, and specifications in the frozen contract. \\
160 unique targets & Distinct mapped declarations; not all are theorem declarations. \\
171 target occurrences & Targets counted again when reused by another contract row. \\
1,034 namespace entries & Entries reported by the historical compiled namespace scan, including generated auxiliaries. \\
\bottomrule
\end{tabular}

A name count is not a measure of mathematical depth. A short wrapper can collect several substantial conclusions; a large structure can expose a difficult hypothesis without constructing its witness. The useful unit is a complete type together with its proof, its dependencies, and an accurate comparison with the prose statement.

\subsection{Three useful reading routes}

For a first reading, follow Sections~\ref{sec:types}--\ref{sec:defects}: exact coordinates lead to weighted transfer, then arithmetic and typed defects. For a coverage review, use the final-paper crosswalk in Appendix~\ref{app:crosswalk}, read the corresponding row in Section~\ref{sec:map}, and inspect its source type. For a reproduction attempt, read Section~\ref{sec:replay} before running any build command; the source checksum check and the compiler audit establish different things.

The main walkthroughs use exact source excerpts, with explanatory equations outside the listings. Names are case-sensitive. \code{PowerSeries}, \code{Polynomial}, \code{LinearMap}, and \code{ContinuousLinearMap} are different types. Similar notation is not a license to replace one with another in a theorem application.

\section{Lean types as mathematical bookkeeping}\label{sec:types}

The basic theorem syntax places assumptions before a colon and a conclusion after it. A term after \code{:=} supplies the definition or proof. A \code{by} block is tactic syntax that elaborates a proof term. A successful elaboration still leaves a semantic question for the reader: does this exact type express the intended mathematical assertion?

\begin{longtable}{p{.30\textwidth}p{.61\textwidth}}
\toprule Expression & Reading in this component \\\midrule\endhead
\code{[Semiring R]} & A typeclass instance supplying semiring operations and laws. Subtraction need not exist. \\
\code{[Ring R]} & A ring, possibly noncommutative; multiplicative order remains significant. \\
\code{[CommSemiring R]} & Commutative multiplication is available, as in the general Fourier certificate. \\
\code{[Module R M]} & A scalar action with its laws; the scalar type is explicitly named. \\
\code{M →ₗ[R] N} & An $R$-linear map from $M$ to $N$, including additive and scalar compatibility. \\
\code{Module.End R M} & The endomorphism ring or semiring of a single module. \\
\code{X →L[k] Y} & A continuous linear map with operator norm in the relevant normed setting. \\
\code{M ≃ₗ[R] N} & A linear equivalence, including both inverse laws. \\
\code{R →+* S} & A ring or semiring homomorphism; it preserves the operations used to transport an inverse. \\
\code{[Fintype E]} & A finite enumeration of all elements of $E$. \\
\code{Fin n} & Natural numbers strictly below $n$, bundled with their bound. \\
\code{ZMod n} & Integers modulo $n$; the nonzero modulus hypotheses matter for finite cardinality. \\
\code{AddSubgroup M} & An additive subgroup, without automatically being a valuation-ring lattice. \\
\code{Nonempty A} & Existence of a term of type $A$; it does not by itself supply a computational search procedure. \\
\code{Prop} & The sort of propositions, used for certificates whose fields are proof obligations. \\
\bottomrule
\end{longtable}

Parenthesized arguments are normally supplied explicitly. Braced arguments can be inferred; square-bracketed arguments request typeclass synthesis. Section variables become parameters of later declarations when used. Therefore a source theorem must be read with its section context. In this project, \code{autoImplicit = false} and \code{relaxedAutoImplicit = false} prevent undeclared identifiers from silently becoming new implicit parameters.

The common \code{constructor} tactic splits a structure or conjunction into its fields. \code{rcases} opens a supplied witness into its data and proofs. \code{rw} rewrites using proved equalities. \code{simp} applies simplification rules and produces a proof term. \code{ext} reduces equality of structured maps to pointwise equality. \code{exact} supplies a term of the required type. \code{rfl} uses definitional equality. None of these words is a substitute for a missing mathematical hypothesis.

\subsection{Definitions, certificates, and algorithms}

A structure such as \code{TwoSidedInverse x y} records equations about already supplied $x$ and $y$. A theorem that accepts this structure does not discover $y$. By contrast, the finite-graph rationality branch defines a denominator and numerator from the graph and proves their properties without accepting rationality as an input.

\code{noncomputable} does not mean unproved. It marks a definition that is not offered as executable code through Lean's ordinary computational evaluation. This component uses noncomputable finite-set choices, finite-type instances, determinants, adjugates, and least-exponent constructions. Their mathematical properties can be proved even though a standalone executable implementation is not supplied.

The three algorithm rows certify mathematical input-output contracts. They do not certify the companion Python package, its storage layout, its numerical residuals, or its performance. Even where a finite enumeration is executable in principle, the contract contains no general complexity estimate or extraction theorem.

\section{Source architecture and entry points}\label{sec:architecture}

The main entry point is \code{import KomorebiVerification.PaperII}. It imports the eight substantive modules but not the audit command. The top-level \code{KomorebiVerification.lean} imports the aggregate for Lake's default library target. The standalone audit imports the substantive modules explicitly and inspects the resulting environment.

\begin{longtable}{p{.30\textwidth}p{.61\textwidth}}
\toprule Module or layer & Responsibility \\\midrule\endhead
Shared foundation & \code{ProjectorCore}, \code{ExactCoordinates}, \code{FibreRefinement}, and \code{ResidueProjectors}: complete finite projector families, range coordinates, masks, and refinement. \\
\code{NestedRefinement} & Parent maps, absorption, child uniqueness, exact coordinate equivalences, cocycles, naturality, and block aggregation. \\
\code{ResidueFourierCRT} & Division-free masks, finite Fourier certificates, primitive-root specialization, polyphase refinement, and CRT labels. \\
\code{WeightedTransfer} & Ordered geometric inverses, normed Neumann inverses, multiplicative filtrations, and precise word pruning. \\
\code{TransferZeta} & Finite weighted walks, exact enumeration, first-edge recurrence, determinant rationality, and coefficient recurrence. \\
\code{ResultantSeparation} & Evaluated B\'ezout identities, numerator and kernel actions, denominator clearing, and certificate-based precision delay. \\
\code{SmithDyadic} & Concrete modular fibres, abstract certified products, rational sign matrices, local denominator weights, and the Smith-phase formula. \\
\code{TypedLinearDefects} & Fine-to-coarse maps between distinct spaces, power and resolvent defects, polynomial compatibility, and operator-norm estimates. \\
\code{LinearDefects} & Ring-valued counterparts and auxiliary inverse identities. \\
\code{KernelAudit} & Representative axiom reports and the fail-closed namespace gate. \\
\bottomrule
\end{longtable}

The main mathematical dependencies are modular. Coordinate refinement does not require a valuation. Weighted pruning needs a filtration but not a graph incidence law. Transfer-zeta counts graph walks but uses integer formal power series, not operator norms. Smith fibres work in finite modular arithmetic. The typed defect module is independent of a particular packet construction: a caller may use a coarsening map, but the theorem accepts any map of the stated type.

\section{Exact packets and refinement}\label{sec:refinement}

Rows 1--5 start from finite complete projector families. A family $P=(P_i)$ satisfies $P_i^2=P_i$, $P_iP_j=0$ for $i\ne j$, and $\sum_iP_i=I$. Orthogonality here is algebraic annihilation. An inner product, self-adjointness, and topology are absent from these assumptions. The canonical formal family also permits zero projectors, whereas the expanded traditional packet terminology requires nonzero members.

The packet space is the product of actual ranges,
\[
\mathcal P(P)=\prod_i\operatorname{range}(P_i).
\]
Its elements carry their range-membership proofs. Analysis sends $x$ to $(P_i x)_i$; synthesis sums the underlying vectors. Completeness proves synthesis after analysis. Pairwise annihilation and idempotence prove analysis after synthesis. Both maps are independently defined, so the second inverse law has mathematical content: every tuple in the product of ranges is accepted.

\subsection{What a refinement witness supplies}

A parent map $\pi:J\to I$ assigns fine labels to coarse labels. The finite set \code{children parent a} is the fibre $\pi^{-1}(a)$. The refinement witness records surjectivity and exact equality of coarse maps with sums of their fine children.
\excerpt{PaperII/NestedRefinement}{refinement}

The structure does not construct a parent map between arbitrary families. To apply its theorems, a caller supplies a map and proves these two fields. The resulting absorption identities are genuine consequences. If $\pi(j)=i$, then
\[
P_iQ_j=Q_j=Q_jP_i;
\]
if $\pi(j)\ne i$, both products vanish. Expand $P_i$ as the sum of its children, distribute composition, and use fine-family orthogonality. Transitivity follows by applying this reasoning through two successive parent maps.

\subsection{A proof walkthrough: recovering a selected finite decomposition}

The most general uniqueness theorem in row 2 does not use the refinement witness. Suppose a selected finite set $S$ of packet vectors sums to $x$, with $v_j\in\operatorname{range}(P_j)$ for $j\in S$. Applying $P_i$ kills all terms except $v_i$, so $v_i=P_i x$. The exact proof shows which hypotheses enter.
\excerpt{PaperII/NestedRefinement}{packet_uniqueness}

First \code{rcases hrange b hb with} extracts a preimage $y$ for the selected packet. Rewriting the packet as $P_b y$ makes idempotence applicable. Next \code{map_sum} distributes $P_b$ across the finite sum. \code{Finset.sum_eq_single_of_mem} leaves the diagonal term; each other term vanishes after its own range witness is opened. Finally \code{congrArg} applies $P_b$ to the supplied sum equality.

The conclusion is restricted to $b\in S$. No conclusion is drawn about values of the packet function outside that finite set. The source family is complete as a bundled object, but the proof of this result never uses its completeness field. This explains the map's ``Stronger'' classification without suggesting a theorem about arbitrary infinite sums.

\subsection{Coordinate changes need no parent map}

Let $\mathsf A_P$ and $\mathsf S_P$ denote analysis and synthesis. For any two complete finite families on the same module, define
\[
R_{P\to Q}=\mathsf A_Q\mathsf S_P,
\qquad R_{Q\to P}=\mathsf A_P\mathsf S_Q.
\]
The source calls these maps refinement and coarsening.
\excerpt{PaperII/NestedRefinement}{coordinate_maps}

These maps are inverses because $\mathsf S_Q\mathsf A_Q=I$ and $\mathsf A_P\mathsf S_P=I$. The proof uses the earlier exact-coordinate simplification rules.
\excerpt{PaperII/NestedRefinement}{coordinate_inverses}

The short \code{simp} proofs rest on established inverse laws; they are not assumptions that refinement loses no data. The two packet products are full exact presentations of the same module. Counting the number of labels alone does not measure their dimension: every label contains an entire projector range.

A parent witness adds component formulas. The fine component is obtained by applying the fine projector to its coarse parent packet; the coarse component is the sum of its fine children. These are additional identities about the already defined coordinate equivalence. They do not turn it into averaging, sampling, or truncation.

With a third complete family $H$, the same inverse laws give the cocycle
\[
R_{Q\to H}R_{P\to Q}=R_{P\to H}.
\]
The theorem \code{refineCoordinates_comp} first proves the equality on each packet and then uses linear-map extensionality. A reader can thus separate the simple conceptual proof from the bookkeeping needed to state it as equality of bundled maps.

\subsection{Naturality and independently refined blocks}

For $T:M\to N$ with source family $P$ and target family $Q$, the exact coordinate operator is $\mathsf A_QT\mathsf S_P$. Its composition law cancels the middle synthesis-analysis pair. Changing the source and target coordinate systems independently gives the corresponding commuting naturality identity.

The output-first block convention is
\[
T_{ji}=Q_jTP_i.
\]
If source and target families refine separately, a coarse block is the double sum over fine descendants of its source and target labels. The mapped theorem permits different source and target modules, which is stronger than the endomorphism special case. Its composite traceability wrapper also records reconstruction of $T$ from all coarse blocks and from all fine blocks.

The finite algebraic result does not establish sparsity of a computed representation, a finite-complexity block extraction routine, or convergence of an infinite block matrix. These need separate hypotheses and, where intended, separate formal developments.

\section{Residue selectors, Fourier certificates, and CRT}\label{sec:residue}

Rows 6--9 act on coefficient functions $x:\mathbb N\to A$. At a positive modulus $N$, the residue projector keeps $x(n)$ when $n\bmod N=j$ and returns zero otherwise. Since exactly one residue label matches each $n$, idempotence, pairwise annihilation, and completeness follow pointwise. No division or roots of unity are required.

This representation is a sequence module, not a formal power-series object with a multiplication law. It is sufficient for coefficient masks. A user wishing to apply the result to a particular function space must supply the connection to that space and any closure or convergence facts.

\subsection{The finite Fourier table is an input over general semirings}

A Fourier realization uses scalars $\omega$, $\omega^\vee$, and $\eta$. The weighted residue kernel is
\[
\eta\sum_{k=0}^{N-1}(\omega^\vee)^{jk}\omega^{rk}.
\]
Over a general commutative semiring the source does not assume that $\omega^\vee$ is an inverse or that $\eta=1/N$. Instead it accepts a finite orthogonality table with exactly the required conclusion at residue addresses.
\excerpt{PaperII/ResidueFourierCRT}{fourier_certificate}

The substantive bridge is from that finite table to every coefficient address. The equation $\omega^N=1$ reduces the exponent $n$ to $n\bmod N$. The last line below uses \code{Fin.ext_iff} to equate residue labels through their underlying natural numbers.
\excerpt{PaperII/ResidueFourierCRT}{fourier_bridge}

After this bridge, the average and the mask are equal coefficient by coefficient: the normalized scalar is one at the selected address and zero elsewhere. The proof works with a finite table rather than an infinite family of independent premises.

There is also a stronger field branch. A primitive root of positive order supplies the normalized inverse-root orthogonality certificate internally. Its field hypothesis is essential to that branch's scope. Root order by itself is not a replacement for the table over arbitrary rings with zero divisors.

\subsection{Refining residues and factoring their labels}

For positive $N$ and $k$, the modulus-$N$ selector is the sum of all modulus-$Nk$ selectors whose labels reduce to it. This is exact regrouping of masks. No packet is discarded and there is no approximation error to estimate.

For positive coprime $m,n$, CRT gives an equivalence between labels modulo $mn$ and pairs of labels modulo $m$ and $n$. The corresponding selector is the composite of the two independent selectors in either order. For $m=2^a$ and positive odd $n$, the dyadic-odd factorization is a specialization.

These are coordinate facts. They do not prove that rational spectral projectors have integral denominators, that a completed lattice splits, or that solving a lifting problem is possible. The later arithmetic modules use different data for those questions. This separation is useful when reading the same dyadic-odd terminology in the traditional paper: residue factorization and spectral denominator clearing are distinct statements.

\section{Weighted inversion: three separate branches}\label{sec:weighted}

Rows 10--11 concern $D+H$ in a possibly noncommutative ring. A supplied two-sided inverse $D^{-1}$ permits the factorization
\[
D+H=D(1+A),\qquad A=D^{-1}H.
\]
Thus an inverse of $1+A$ is multiplied by $D^{-1}$ on the right. Reversing this order changes the formula in a noncommutative ring. The source parameter named \code{K} in these declarations means the perturbation $H$; it is not a fraction field.

Three hypotheses lead to three different results. Nilpotence gives a finite ring identity. A filtration-killing homomorphism gives a finite identity after passing to precision. A strict norm bound in a complete normed ring gives an infinite Neumann sum. The formal component does not merge these into an unrestricted adic convergence theorem.

\subsection{The finite geometric identity}

For $A^N=0$, define $B=\sum_{n=0}^{N-1}(-A)^n$. The product formulas for finite geometric sums prove both $(1+A)B=1$ and $B(1+A)=1$.
\excerpt{PaperII/WeightedTransfer}{finite_inverse}

The source separately invokes the two ordered geometric-sum lemmas. This is why \code{TwoSidedInverse} has two fields. A one-sided inverse does not automatically supply the other in a general noncommutative ring. Once this result is available, the diagonal-transfer theorem transports the candidate through the factorization of $D+H$.

The cutoff $N$ and its nilpotence proof are inputs. The theorem is not a procedure for deciding whether an arbitrary ring element is nilpotent. Its finite sum is explicit, but a concrete matrix application must still identify the matrix ring, supply $D^{-1}$, and prove the required equations.

\subsection{The filtration is deliberately minimal}

A multiplicative filtration records additive closure, decreasing levels, and multiplicative addition of indices.
\excerpt{PaperII/WeightedTransfer}{filtration}

Write $F^n$ for the predicate \code{Mem n}. Multiplicativity is $F^mF^n\subseteq F^{m+n}$; antitonicity sends membership at a larger index to membership at a smaller one. Completeness, separation, and an inverse-limit topology are not fields of this structure. Thus the word ``filtration'' in its name does not imply any convergence statement.

A precision certificate for a ring homomorphism $q:R\to S$ says $q(F^L)=0$. It does not claim that $S$ is a particular quotient or that $\ker q=F^L$. Killing the level is sufficient for the theorem. The proof below first puts a word into its cumulative filtration level and then uses antitonicity to reach $F^L$.
\excerpt{PaperII/WeightedTransfer}{precision}

The caller supplies the edge-membership facts $A_e\in F^{w(e)}$. Induction on a list gives $A_{e_1}\cdots A_{e_r}\in F^{w(e_1)+\cdots+w(e_r)}$. All words of weight at least $L$ therefore disappear after applying $q$.

\subsection{Written products and chronological paths}

The definition \code{wordProduct edge word} is \code{(word.map edge).prod}. For the list $[e_1,\ldots,e_r]$, it gives the written product $A_{e_1}\cdots A_{e_r}$. When these are operators, the rightmost factor acts first. This is the convention in the canonical formal weighted-word branch.

The expanded traditional block-path discussion writes a chronologically traversed path $e_1,\ldots,e_r$ as $A_{e_r}\cdots A_{e_1}$. To translate an indexed term, reverse the list. The full sums over all words of a fixed length agree under this bijection; the weight-truncated sums also agree because reversal preserves cumulative weight. Individual indexed products need not agree. An equation for one labelled path must therefore retain its order convention.

The generic formal weighted-word sum ranges over edge lists without imposing graph adjacency. In a concrete block realization, incompatible products may vanish for structural reasons, but those incidence facts are not fields of \code{MultiplicativeFiltration}. They must be supplied when interpreting the sum as a sum over directed block paths.

\subsection{What is actually retained}

The pruned expression contains only words with total weight strictly below $L$. Its length cutoff $N$ is a separate parameter.
\excerpt{PaperII/WeightedTransfer}{pruned_sum}

If each edge weight is at least $m$, a length-$N$ word has weight at least $mN$. Hence $L\le mN$ makes the $N$th mapped power vanish. For $m>0$, the proved ceiling inequality allows $N=\lceil L/m\rceil$. With merely positive integral weights, $N=L$ is a valid coarse cutoff.

The equality between the mapped full geometric sum and the mapped retained-word sum is proved by splitting the finite word sums into retained and killed terms. The inverse conclusion then follows from mapped nilpotence. This order of reasoning matters: pruning is justified by exact zero in the precision realization, not by a small numerical residual.

The full diagonal-plus-transfer theorem exposes all remaining data. It accepts the inverse of $D$, the normalized edge expansion, the precision certificate, the minimum-weight bound, the cutoff inequality, and the edge-membership proofs.
\excerpt{PaperII/WeightedTransfer}{diagonal_inverse}

The proof first obtains the pruned inverse of $1+\sum_eA_e$. It maps the supplied normalized-transfer equality, rewrites the inverse target, and applies the diagonal factorization lemma. The final \code{simpa} only unfolds the candidate and the homomorphism's preservation of addition and multiplication. No inverse of $D$ or quotient construction is hidden in that final simplification.

\subsection{The analytic branch and the algorithm contract}

In a complete normed ring, $\|D^{-1}H\|<1$ supplies a separate Neumann-series inverse. The source defines the analytic candidate noncomputably and proves both inverse laws using the normed summability theory. This is neither the finite nilpotence premise nor completeness of an arbitrary filtered ring. A valuation-based application needs a bridge to one of the formal settings.

Row 24 packages the bounded geometric/pruned expression as a mathematical algorithm specification. It proves that the chosen ceiling cutoff meets the required inequality and that the mapped result is a two-sided inverse. It does not construct a basis-level modular inverse routine, decide symbolic zero, certify the Python implementation, or establish a runtime bound.

\section{Transfer-zeta: from independent walks to rationality}\label{sec:zeta}

Rows 12--13 and 25 are among the component's strongest construction results. The input is a finite directed multigraph with positive natural edge weights. Parallel edges are distinct edge values and therefore contribute their multiplicities automatically. Repeated vertices and edges are allowed: the objects counted are walks.
\excerpt{PaperII/TransferZeta}{graph}

The recursion \code{IsWalk} says an empty edge list starts and ends at the same vertex. A nonempty list starts with an edge whose source is the current vertex, then continues from that edge's target. The total weight is the sum of edge weights.

\subsection{Why a coefficient has a finite enumeration}

Positivity implies that a walk of total weight $n$ has length at most $n$. The enumerator can therefore range over lengths $0,\ldots,n$ and all edge tuples of each length, then filter by endpoint validity and exact total weight.
\excerpt{PaperII/TransferZeta}{enumerator}

The module also defines an independent semantic type of finite edge lists having the requested endpoints and weight, without building the enumeration bound into that definition. The capture theorem uses positivity to connect the unrestricted semantic representation to the bounded tuple enumeration. The counting theorem and the trace wrapper identify the returned natural number with the cardinality of this semantic walk type.

This avoids defining correctness merely as equality with the same algorithm under another name. The finite enumeration and the semantic list predicate are separate constructions connected by equivalences. The implementation remains a potentially large brute-force enumeration; no efficiency theorem is claimed.

\subsection{The first-edge recurrence}

Fix a target $t$ and let $a_{st}(n)$ count walks from $s$ to $t$ of weight $n$. The empty walk contributes $1$ exactly when $s=t$ and $n=0$. Every other walk has a unique first edge. Removing it gives
\[
a_{st}(n)=\mathbf 1_{s=t,\,n=0}
+\sum_{\substack{e:\,s\to u\\w(e)\le n}}a_{ut}(n-w(e)).
\]
The source constructs the corresponding first-edge equivalence and obtains the recurrence by cardinality. Positivity ensures the recursive weight decreases. The bound $w(e)\le n$ is explicit, avoiding an incorrect use of truncated natural subtraction when the first edge is too heavy.

Define $Z_{st}(X)=\sum_{n\ge0}a_{st}(n)X^n$. This is a formal power series over $\mathbb Z$. The weighted adjacency entry $A_{su}(X)$ sums $X^{w(e)}$ over all edges $s\to u$. The first-edge recurrence becomes
\[
(I-A(X))Z(X)=I.
\]
The row/column orientation here follows first-edge traversal. When comparing it with output-first operator blocks, read the definitions rather than transferring an index convention by analogy.

\subsection{A second construction: determinant and adjugate}

Finite vertex data define the polynomial matrix $B(X)=I-A(X)$ and
\[
Q(X)=\det B(X),\qquad P_{st}(X)=\operatorname{adj}(B(X))_{st}.
\]
Positive edge weights give $A(0)=0$, hence $B(0)=I$ and $Q(0)=1$. Thus $Q$ is a unit in $\mathbb Z[\![X]\!]$. The adjugate identity constructs a two-sided inverse matrix $Z^{\mathrm{mat}}$.

The crucial proof is the identification of the independently enumerated series matrix with this rational matrix. It uses associativity and the two inverse equations.
\excerpt{PaperII/TransferZeta}{identify_series}

This is the familiar uniqueness calculation $Z=I Z=(Z^{\mathrm{mat}}B)Z=Z^{\mathrm{mat}}(BZ)=Z^{\mathrm{mat}}$. It makes the construction claim auditable: rationality was not an assumption about the enumerated coefficients. The final existential theorem simply packages the already constructed numerator and denominator.
\excerpt{PaperII/TransferZeta}{rationality}

The older \code{RationalCertificate} interface remains available for an externally supplied pair $P,Q$. Its theorem is certificate-explicit. The canonical finite-graph branch does not take that certificate; it derives its own data from the graph. The map preserves both interfaces so users can choose one without confusing their hypothesis burdens.

\subsection{The recurrence is extracted from the proved identity}

Write $Q(X)=\sum q_jX^j$. Taking coefficient $n$ in $QZ=P$ gives an antidiagonal convolution. Once $n>\deg P$, the right-hand coefficient is zero.
\excerpt{PaperII/TransferZeta}{recurrence}

Because $q_0=1$, this can be read as an eventual constant-coefficient recurrence for $a_{st}(n)$. The source's antidiagonal form states the identity without needing to rearrange its leading term. It does not guess a recurrence from sampled coefficients.

The determinant denominator may have canceling factors for a particular matrix entry, so this construction does not prove minimality of $Q$ or of the recurrence order. The determinant and adjugate definitions are noncomputable here. The result concerns formal power series and supplies no disk of analytic convergence, extracted symbolic determinant program, or infinite-graph rationality theorem.

\section{Evaluated separation and lattice precision}\label{sec:separation}

Rows 14--16 start with endomorphisms $F,G,U,V$ of a vector space over a field $K$. The central premise is an evaluated B\'ezout equality
\[
UF+VG=rI.
\]
The source does not construct polynomials whose evaluation gives these maps, compute a general resultant, factor a characteristic polynomial, or instantiate Cayley--Hamilton. These operations may justify the premise in a mathematical application, but they remain upstream of this formal interface.

Let $N=VG$. The denominator-free identity is $N^2=rN$. After adding $r\ne0$, the normalized map $E=r^{-1}N$ is idempotent. The bundled \code{EvaluatedBezoutCertificate} contains six pairwise commutations as well as nonvanishing, B\'ezout, and annihilation. The mapped minimal theorems need fewer fields and state this directly.

\subsection{A proof walkthrough: quasi-idempotence without division}

Only $FG=GF$, $GU=UG$, and $FG=0$ are needed to kill the cross term $VGUF$. Using $VG=rI-UF$ on one occurrence of the numerator then proves the square identity.
\excerpt{PaperII/ResultantSeparation}{bezout_minimal}

The proof rewrites just the second factor, distributes multiplication over subtraction, and invokes cross-term vanishing. It needs neither $r\ne0$ nor four unused commutations from the stronger bundle. Consequently the denominator-free theorem remains meaningful even when normalization is unavailable. The later normalized theorem uses scalar associativity and field simplification with the explicit nonzero hypothesis.

On $\ker F$, the evaluated B\'ezout equation gives $Nx=rx$ directly. On $\ker G$, $Nx=V(Gx)=0$. After normalization, $E$ acts as identity and zero respectively. A statement that $F$ is nilpotent on some larger subspace would not imply $Fx=0$ there; the formal theorem uses exact kernel membership. This prevents a primary-summand intuition from silently strengthening the proved conclusion.

\subsection{What the formal lattice means}

The lattice parameter is an \code{AddSubgroup M}. Preservation by a map means that each member is sent back into this additive subgroup. A general valuation ring, finite generation, full rank, completeness, and a chosen basis are not built into its type.

A clearing certificate at exponent $s$ states that $\pi^sE$ preserves the lattice. The source names this \code{IsSeparationWeight}. If multiplication by $\pi$ also preserves the lattice, the set of clearing exponents is upward closed. Given existence of some clearing exponent, \code{Nat.find} defines a least one and proves its minimality property. This is a noncomputable least-witness construction, not a search algorithm over matrix entries.

If $r=u\pi^s$, then
\[
\pi^sE=u^{-1}VG.
\]
To turn this algebraic identity into lattice preservation, one must know that $u^{-1}$ preserves the chosen lattice, together with $V$ and $G$. An abstract field unit need not preserve an arbitrary additive subgroup. The formal lattice-unit certificate makes this assumption visible. The resulting exponent bounds the least separation weight from above; it need not be the minimal exponent.

The expanded traditional theory identifies exact weights through a lattice gluing module and gives necessary as well as sufficient delay statements. Those ordinary written proofs are beyond these mapped clearing-certificate consequences. They should not be attributed to \code{separationWeight} merely because the names describe related mathematics.

\subsection{Precision delay as transport of a witness}

Formal congruence at depth $n$ means an explicit equation $x-y=\pi^nz$ with $z$ in the lattice. If $\pi^sE$ preserves the lattice, input congruence at $m+s$ becomes output congruence at $m$.
\excerpt{PaperII/ResultantSeparation}{precision_delay}

The proof opens the original witness $z$. It chooses the new witness $z'=(\pi^sE)z$, whose lattice membership is exactly the clearing hypothesis. Linearity and $\pi^{m+s}=\pi^m\pi^s$ establish the new difference equation. No digit expansion, topology, or valuation API enters this proof. It is sufficient delay from a supplied certificate, without a claim that fewer input digits must fail.

\section{Concrete Smith fibres and the dyadic sign model}\label{sec:smith}

Row 18 proves a concrete arithmetic statement over \code{ZMod (p^depth)} for $p>1$. The hypothesis is not that $p$ is prime. Since the modulus and coefficient are powers of the same integer, the needed gcd and divisibility identities hold at this stated scope.

\subsection{The solution fibre is an actual type of solutions}

For a weight $s$ and depth $L$, multiplication by $p^s$ is an additive endomorphism of $\mathbb Z/p^L\mathbb Z$. Its fibre over $b$ is the subtype of elements that map to $b$.
\excerpt{PaperII/SmithDyadic}{smith_fibre}

The criterion is
\[
p^{\min(s,L)}\mid\operatorname{val}(b),
\]
where \code{target.val} is the canonical natural representative. The source proves that this is equivalent to nonemptiness of the actual fibre. If $s\le L$, an admissibility witness supplies a quotient of the target representative and hence a solution. If $L\le s$, admissibility forces the target to be zero; the zero residue supplies a solution.

The forward direction uses the modular representative of a supplied solution and divisibility by both the coefficient and modulus. Thus solvability is proved from concrete modular arithmetic, rather than inserted as a certificate field.

\subsection{A proof walkthrough: counting an admissible fibre}

Once a solution $x_0$ exists, translation by $x_0$ identifies the fibre with the kernel of multiplication by $p^s$. The kernel cardinality in this finite cyclic additive group is the gcd of its order with $p^s$.
\excerpt{PaperII/SmithDyadic}{smith_count}

The final two cases use the divisibility of one power of $p$ by the other to obtain
\[
\#\{x:p^sx=b\pmod {p^L}\}=p^{\min(s,L)}.
\]
This is an exact cardinality of a finite type. It does not enumerate the solutions or provide a particular coordinate normal form for every one of them.

For a finite diagonal system, the global fibre is a dependent function type: one solution in each coordinate. Products of finite cardinalities give the count, while one inadmissible coordinate makes the product type empty.
\excerpt{PaperII/SmithDyadic}{smith_global}

If all coordinates are admissible, the cardinality is $p^{\sum_i\min(s_i,L)}$; otherwise it is zero. At depth zero, the quotient has one element and the exponent is zero, consistent with the stated formula. A literal zero diagonal coefficient can be represented at precision $L$ by any weight at least $L$, since $p^s$ then vanishes in the quotient.

\subsection{The abstract extension keeps its certificate boundary}

The same module has separate \code{SmithSolvabilityCertificate} and \code{SmithCardinalityCertificate} structures. For arbitrary finite coordinate fibres, these supply the per-coordinate equivalence and count, with an externally chosen base $q$. The global product theorems then assemble them.

In a general discrete valuation ring with residue field of cardinality $q$, the familiar coordinate count has the analogous form. The formal component does not construct that DVR arithmetic or derive its cardinality theorem. Likewise it does not calculate Smith normal form of an input matrix, find change-of-basis matrices, or extract diagonal weights. The automatic row-18 result begins with a diagonal system over the concrete modular quotient; the abstract extension begins with certificates.

\subsection{Rational sign projectors and local denominator weight}

Row 17 fixes the rational matrices
\[
J=\begin{pmatrix}1&1\\0&-1\end{pmatrix},\qquad
P_+=\begin{pmatrix}1&1/2\\0&0\end{pmatrix},\qquad
P_-=\begin{pmatrix}0&-1/2\\0&1\end{pmatrix}.
\]
Entrywise proofs establish $J^2=I$, projector idempotence, mutual annihilation, and $P_++P_-=I$. The local integrality predicate says that a rational entry's reduced denominator is not divisible by $p$. Each projector has exact local denominator weight zero at every prime other than two and exact weight one at two.

The local word matters. Global integrality demands integer entries; $p$-integrality permits denominators prime to $p$. A global denominator weight and a local denominator weight are defined separately in the source. The sign-factor resultant is explicitly computed as $\operatorname{Res}(X-1,X+1)=2$, and its dyadic valuation is matched to the exact local weights of these matrices.

This is a rational-entry model using natural-number valuations. It does not construct $\mathbb Q_p$ or its valuation ring, identify arbitrary completed lattices, or prove a general equivalence between all rational integrality criteria and lattice preservation. Reduction modulo two applies to the characteristic factors; the rational matrices containing $1/2$ do not have an entrywise reduction modulo two.

\subsection{The Smith-phase identity has only algebraic scope}

Row 22 checks the pulled-back fractional expression \code{smithPhase} at phase $-1$. Over a field with $2\ne0$ and $z\ne0$, its value is $z^{-1}$. The extra manuscript condition $z\ne-1$ is unnecessary for this algebraic identity and is removed by the formal theorem.

The result does not formalize the complex Smith disk, a M\"obius-domain interpretation, radius preservation, or physical impedance. In the traditional companion, navigate to Section 12.3, the unnumbered Smith-phase identity on printed page 45. The subsection counter is not an equation number.

\section{Typed defects: domain and codomain stay visible}\label{sec:defects}

Rows 19--21 use a genuinely typed setting. Let $X_f$ and $X_c$ be fine and coarse modules, $C:X_f\to X_c$, and let $T_f$ and $T_c$ act on their respective spaces. The one-step defect is
\[
\Delta=CT_f-T_cC:X_f\longrightarrow X_c.
\]
The source does not assume that the two spaces coincide.
\excerpt{PaperII/TypedLinearDefects}{typed_defect}

Notice that the scalar type is a semiring while the spaces have additive group structures, allowing subtraction of maps. Scalar subtraction is not needed for the basic power identity. This is one example of why the complete type is more informative than a loose description such as ``linear algebra over a ring.''

\subsection{The telescoping mechanism}

Define $E_n=CT_f^n-T_c^nC$. The induction uses
\[
E_{n+1}=T_cE_n+\Delta T_f^n.
\]
Expanding $E_n$ gives
\[
CT_f^n-T_c^nC
=\sum_{k=0}^{n-1}T_c^{n-1-k}\Delta T_f^k.
\]
Every summand is a map $X_f\to X_c$: the right power acts in $X_f$, the defect crosses to $X_c$, and the left power acts there. Treating $C$ as an endomorphism would hide this type constraint.

The all-power theorem includes $n=0$. The left side is $C-C=0$ and the finite sum is empty. The positive-power theorem is a wrapper around this stronger statement. In the successor proof, \code{Finset.sum_range_succ} separates the new last term; natural-number arithmetic verifies the shifted exponents. The exponent $n-1-k$ occurs only at $k<n$, where it has the intended value.

If the one-step defect is zero, every term in the sum vanishes.
\excerpt{PaperII/TypedLinearDefects}{zero_defect}

This extends to finite polynomial functional calculus by summing the corresponding power identities with scalar coefficients. The polynomial-action definition explicitly carries the commutative scalar assumptions needed for those coefficients to act centrally. The result does not imply a holomorphic, exponential, or other infinite functional-calculus theorem without further analytic work.

\subsection{The operator-norm estimate}

In the continuous-linear-map branch, the same finite identity yields
\[
\|CT_f^n-T_c^nC\|
\le \|\Delta\|\sum_{k=0}^{n-1}
\|T_c\|^{n-1-k}\|T_f\|^k.
\]
Triangle inequality bounds the norm of the sum; submultiplicativity bounds each composition; the endomorphism power estimate bounds each power. Completeness of the spaces is unnecessary because this is a finite sum.

If both operator norms are at most $\rho$, the component derives the common-bound majorant
\[
\|CT_f^n-T_c^nC\|\le n\rho^{n-1}\|\Delta\|.
\]
The source states the natural-number endpoint carefully, so the all-$n$ theorem also includes zero. The hypotheses supply $\rho$ and the two norm bounds. A spectral-radius bound is not interchangeable with an operator-norm bound here, and the majorant is not asserted to be sharp.

For example, exact intertwining $\Delta=0$ forces every finite-power defect to vanish. A small nonzero defect may still be amplified by large powers. The theorem gives an upper bound with the stated norm data; it does not identify the actual error for a particular trajectory or nonlinear dynamics.

\subsection{Resolvents from the inverse equations actually used}

Let $R_f$ and $R_c$ be supplied inverse candidates for the scalar shifts. The continuous theorem needs only
\[
(zI-T_f)R_f=I,\qquad R_c(zI-T_c)=I.
\]
The fine witness is a right inverse; the coarse witness is a left inverse. They are sufficient for the ordered sandwich calculation
\[
CR_f-R_cC=R_c\Delta R_f.
\]
\excerpt{PaperII/TypedLinearDefects}{resolvent}

The source first applies a general inverse-sandwich identity. The difference of shifted compositions simplifies to $CT_f-T_cC$ because the two scalar terms cancel. The sign is therefore positive on the right for this convention of $\Delta$. Reversing the defect definition would change that sign.

Applying the operator norm gives
\[
\|CR_f-R_cC\|\le\|R_c\|\,\|\Delta\|\,\|R_f\|.
\]
No theorem here constructs the inverse candidates or proves that $z$ lies in a resolvent set. The algebraic typed branch uses explicit two-sided inverse structures; the continuous branch exposes the weaker one-sided equations above. A user should choose the theorem matching the available data rather than infer inverse existence from the word ``resolvent.''

\subsection{What remains in the written companion}

The final traditional development also proves holomorphic functional-calculus, exponential, and multilevel defect identities with their own analytic or compositional hypotheses. The preserved Lean map records finite polynomial compatibility, the power estimates, and the inverse-witness resolvent identities. Those additional written results remain ordinary proofs in this edition. The formal contract can be complete while the broader traditional theorem inventory has a different coverage boundary.

\section{Reproducing the formal project}\label{sec:replay}

Reproduction has two separate layers. The lightweight guide checks confirm that the supplied source matches its checksum manifest, that all mapped names can be found lexically, and that the printed excerpts are exact slices of those files. The Lean replay builds the project, resolves names in the compiled environment, and traverses axiom dependencies. Only the latter executes Lean.

\subsection{Start with the preserved pins}

Work in a writable copy of \code{source/canonical_v1.2.1/formalization/lean/}. Keep \code{lean-toolchain}, \code{lakefile.toml}, and \code{lake-manifest.json} together. The exact toolchain is \code{leanprover/lean4:v4.34.0-rc2}; the Mathlib revision is \code{aeb72f1e9981b6d35b426b16dca7759be3e7ba06}. Appendix~\ref{app:imports} records every dependency revision in the delivered manifest.

A working Lean/Elan installation, network access for dependencies or a suitable local cache, and sufficient disk space are prerequisites. The following commands are the preserved report's replay procedure, not commands executed while producing this guide.
\begin{Verbatim}[fontsize=\small,frame=leftline,rulecolor=\color{teal},framesep=6pt,breaklines=true]
cd source/canonical_v1.2.1/formalization/lean
lake exe cache get
lake build
python3 scripts/verify_theorem_map.py
lake env lean -E warning KomorebiVerification/PaperII/KernelAudit.lean
\end{Verbatim}

The cache command obtains compiled dependency artifacts. The build checks the aggregate library against the pinned environment. The Python map checker parses exactly one versioned theorem map, expands its namespace prefixes, requires rows 1 through 25, and writes a temporary Lean command that checks every mapped name in the compiled environment. The script removes that temporary file afterward. Its environment check is stronger than this guide's lexical source lookup.

The final command separately executes the fail-closed audit. Warning escalation is part of the documented invocation. Preserve the exit status and full output of each command; a report written after the fact should not substitute for those transcripts in a new replay.

\subsection{The audit gate and its trust boundary}

The audit has two parts. \texttt{\#print axioms} displays dependencies for a representative public sample. A subsequent command traverses every compiled declaration in the imported \code{Komorebi} namespace and rejects unsafe entries or dependencies outside the allowed basis.
\excerpt{PaperII/KernelAudit}{audit_gate}

The approved foundations are \code{propext}, \code{Classical.choice}, and \code{Quot.sound}. The gate rejects any other axiom dependency, including \code{sorryAx}, and explicitly fails if no project declaration was inspected. A successful audit is a property of the imported compiled environment at that run. It is not inferred solely from the absence of the word \code{sorry} in source text.

The sample of printed axioms is not the complete traversal. Conversely, the namespace count includes generated declarations and helper entries, so it is not a count of independent mathematical theorems. The map checker and namespace gate also do not decide whether a theorem's type faithfully represents a prose theorem; that requires scope comparison such as the canonical map and this guide.

\subsection{A minimal source-inspection session}

After a successful build, create a scratch Lean file beside the project's aggregate import. The following inspection commands ask Lean to display two representative declarations and the axiom dependencies of the automatic rationality result.
\begin{Verbatim}[fontsize=\small,frame=leftline,rulecolor=\color{teal},framesep=6pt,breaklines=true,breakanywhere=true]
import KomorebiVerification.PaperII
#check Komorebi.PaperII.NestedRefinement.coordinateRefinementEquiv
#check Komorebi.PaperII.SmithDyadic.smithZMod_global_conditional_count
#print axioms Komorebi.PaperII.TransferZeta.PositiveWeightedMultigraph.finite_graph_pathZeta_rationality
\end{Verbatim}

Run it with \code{lake env lean YourInspection.lean}. These are suggested inspection commands, not a new executed transcript. Use \code{\#print} or editor hover to inspect elaborated types, including inferred section parameters and typeclass assumptions. The printed source excerpt alone may omit surrounding binders; the full declaration type is decisive.

When applying a theorem, first fix the intended scalar type and domain/codomain. Then assemble the structural object and prove its fields. For a weighted inverse, this means the filtration, map $q$, precision certificate, diagonal inverse, and edge expansion. For a finite graph, it means finite vertex/edge types, endpoint maps, and positive weights. For a Smith system, it means a diagonal weight function and modular target. This preparation makes failures informative: a missing instance differs from a missing mathematical witness.

\subsection{Common failure modes}

\begin{longtable}{p{.28\textwidth}p{.63\textwidth}}
\toprule Symptom & What to check \\\midrule\endhead
Toolchain mismatch & Read \code{lean-toolchain}; do not silently substitute a stable release for the pinned release candidate. \\
Missing imported module & Confirm dependencies and cache availability, then build from the Lean project root. \\
Unknown mapped name & Run the map checker only after a successful aggregate build; verify source and manifest identity. \\
Unsolved typeclass & Inspect the theorem's complete type and the scalar/module, finite-type, or normed-space assumptions. \\
Missing inverse proof & Supply the actual equation required by the chosen theorem; naming a candidate does not prove it invertible. \\
Unexpected axiom & Preserve the failing declaration and transcript; do not weaken the allowed-axiom gate to make the run pass. \\
Different namespace count & Compare toolchain, dependencies, source hashes, and imports before treating the run as a replay of this release. \\
Successful Python tests & Interpret them as tests of the Python component; they are not Lean kernel evidence. \\
\bottomrule
\end{longtable}

A new successful run should record the source hashes, dependency pins, commands, exit codes, output hashes, and environment. If any formal source changes, the old report no longer supplies execution evidence for that changed source, even if the filenames remain the same.

\section{Evidence retained and evidence newly checked}\label{sec:evidence}

The canonical \code{VERIFICATION_REPORT.md}, dated 3 September 2026, reports successful aggregate build, theorem-map resolution of 160 unique declarations, and a fail-closed scan of 1,034 compiled namespace entries. It records the approved axiom basis and states that no audited declaration reports \code{sorryAx}. This guide inherits those statements as a release report.

The inspected canonical component contains no raw transcript of those Lean commands and no per-command log hash attestation. Consequently this edition does not claim a fresh compiler replay, a fresh environment scan, or a source-to-transcript identity chain that the supplied artifacts do not contain. The report itself is checksummed, so its preservation is verifiable; checksumming a report does not independently establish its execution claims.

The new lightweight check establishes the following facts directly from the delivered files:
\begin{enumerate}
\item All 22 files under the preserved formal component match their entries in the release \code{SHA256SUMS.txt} manifest.
\item The same 22 files match the separately preserved baseline copy in this editing workspace.
\item The normative map has exactly rows 1--25 and 160 distinct expanded declaration names, with 171 occurrences across rows.
\item Every mapped name has a lexical declaration location in the preserved source. This is navigation evidence, not compiled name resolution.
\item The 15 Lean source files contain 5,603 physical lines. A comment- and string-aware scan finds no active token from the documented placeholder/unsafe source policy list.
\item Every exact excerpt matches its recorded inclusive source slice; each source and excerpt has a SHA-256 digest.
\end{enumerate}

The scan policy checks active occurrences of \code{sorry}, \code{admit}, \code{sorryAx}, \code{axiom}, \code{unsafe}, \code{native_decide}, and named dynamic-declaration shortcuts. It strips nested block comments, line comments, and strings while retaining line positions. This is a useful static check, not a replacement for Lean elaboration or the environment's transitive axiom inspection.

The guide's PDF build and visual review are new publication checks. They verify rendering, references, fonts, line overflow, page coverage, and consistency of the guide's data files. They do not add mathematical formalization coverage. The complete generated audit records \code{lean_compiler_run: false} explicitly.

\subsection{A practical rule for reporting success}

Report the strongest evidence actually obtained. After running only the data generator, say that static source identity and map reconciliation passed. After a new successful Lean replay, cite that replay's own logs and environment. When quoting the preserved count of 1,034 entries, attribute it to the historical report. When using an expanded traditional theorem, identify its ordinary written proof unless a precise Lean target with matching hypotheses has been supplied.

This discipline makes the component usable: readers can trust an automatic construction where it exists, supply a certificate where it is required, and locate the written mathematics that remains beyond the formal contract.

\clearpage
\section{The complete 25-row contract}\label{sec:map}
The following map preserves the canonical v1.1.1 classifications, checked scopes, and explicit boundaries. The targets use the namespace abbreviations below. These abbreviations are only typography: the companion JSON expands every name and records its source file and line. A declaration may occur in several rows; there are 171 occurrences and 160 distinct names. Definitions, structures, equivalences, and theorems all count as declaration targets.
\begin{longtable}{p{.12\textwidth}p{.79\textwidth}}
\toprule Prefix & Full namespace \\ \midrule\endhead
\code{PC} & \code{Komorebi.Foundation.LinearProjectorFamily} \\ \addlinespace
\code{NR} & \code{Komorebi.PaperII.NestedRefinement} \\ \addlinespace
\code{R} & \code{Komorebi.Residue} \\ \addlinespace
\code{WT} & \code{Komorebi.PaperII.WeightedTransfer} \\ \addlinespace
\code{TZ} & \code{Komorebi.PaperII.TransferZeta.PositiveWeightedMultigraph} \\ \addlinespace
\code{TZC} & \code{Komorebi.PaperII.TransferZeta.PositiveWeightedMultigraph.RationalCertificate} \\ \addlinespace
\code{RS} & \code{Komorebi.PaperII.ResultantSeparation} \\ \addlinespace
\code{SD} & \code{Komorebi.PaperII.SmithDyadic} \\ \addlinespace
\code{LD} & \code{Komorebi.PaperII.LinearDefects} \\ \addlinespace
\code{TLD} & \code{Komorebi.PaperII.TypedLinearDefects} \\ \addlinespace
\bottomrule\end{longtable}

\Needspace{7\baselineskip}
\subsection{Row 1: Refinement map}
\label{row:1}
\textbf{Classification:} Exact.
\par\smallskip\textbf{Declaration targets}\par
{\small\noindent\code{NR.children}\par}
{\small\noindent\code{NR.mem_children}\par}
{\small\noindent\code{NR.Refines}\par}
\smallskip\textbf{Checked scope.} For finite complete projector families over a semiring module, \code{Refines} records a surjective parent map and equality of every coarse projector with the finite sum over its fine children.
\par\smallskip\textbf{Boundary.} No refinement is constructed from arbitrary projector families; no infinite or topological sum is asserted.
\medskip
\Needspace{7\baselineskip}
\subsection{Row 2: Nested refinement theorem}
\label{row:2}
\textbf{Classification:} Stronger.
\par\smallskip\textbf{Declaration targets}\par
{\small\noindent\code{NR.finitePacket_decomposition_unique}\par}
{\small\noindent\code{NR.finitePacket_decomposition_ext}\par}
{\small\noindent\code{NR.Refines.coarse_comp_fine_eq_self}\par}
{\small\noindent\code{NR.Refines.coarse_comp_fine_eq_zero}\par}
{\small\noindent\code{NR.Refines.fine_comp_coarse_eq_self}\par}
{\small\noindent\code{NR.Refines.fine_comp_coarse_eq_zero}\par}
{\small\noindent\code{NR.Refines.coarse_packet_reconstruction}\par}
{\small\noindent\code{NR.Refines.child_decomposition_unique}\par}
{\small\noindent\code{NR.Refines.trans}\par}
\smallskip\textbf{Checked scope.} Parent/child absorption holds in both composition orders and refinements compose. More generally, any finite sum of vectors lying in mutually orthogonal packet ranges has uniquely determined components, without a refinement or completeness hypothesis.
\par\smallskip\textbf{Boundary.} Uniqueness is only asserted on the selected finite set and requires each component to lie in its stated projector range. All index types and sums are finite.
\medskip
\Needspace{7\baselineskip}
\subsection{Row 3: Exact coordinate refinement}
\label{row:3}
\textbf{Classification:} Stronger.
\par\smallskip\textbf{Declaration targets}\par
{\small\noindent\code{NR.refineCoordinates}\par}
{\small\noindent\code{NR.coarsenCoordinates}\par}
{\small\noindent\code{NR.refineCoordinates_self}\par}
{\small\noindent\code{NR.refineCoordinates_comp}\par}
{\small\noindent\code{NR.coordinateRefinementEquiv_trans}\par}
{\small\noindent\code{NR.Refines.refineCoordinates_component}\par}
{\small\noindent\code{NR.Refines.coarsenCoordinates_component}\par}
\smallskip\textbf{Checked scope.} Analysis after synthesis defines cross-level linear maps for any two finite complete projector families; these maps obey identity and cocycle laws through every intermediate family. A \code{Refines} witness gives the parent/child component formulas.
\par\smallskip\textbf{Boundary.} These are coherent changes of exact coordinates, not truncation or dimensional reduction operators; the cocycle uses exact completeness at each finite level.
\medskip
\Needspace{7\baselineskip}
\subsection{Row 4: Perfect reconstruction across levels}
\label{row:4}
\textbf{Classification:} Stronger.
\par\smallskip\textbf{Declaration targets}\par
{\small\noindent\code{NR.coarsenCoordinates_refineCoordinates}\par}
{\small\noindent\code{NR.refineCoordinates_coarsenCoordinates}\par}
{\small\noindent\code{NR.coordinateRefinementEquiv}\par}
{\small\noindent\code{NR.coordinateOperator}\par}
{\small\noindent\code{NR.coordinateOperator_comp}\par}
{\small\noindent\code{NR.coordinateOperator_naturality}\par}
\smallskip\textbf{Checked scope.} Cross-level refinement and coarsening are mutual inverses over a semiring module. The typed coordinate representation of every linear map preserves identities and composition and commutes with independent exact coordinate changes on source and target.
\par\smallskip\textbf{Boundary.} The inverse laws need completeness of the projector families but no parent map. This is finite algebraic naturality, not convergence for an infinite packet tower.
\medskip
\Needspace{7\baselineskip}
\subsection{Row 5: Transfer-block refinement}
\label{row:5}
\textbf{Classification:} Stronger.
\par\smallskip\textbf{Declaration targets}\par
{\small\noindent\code{NR.transferBlock_refinement_between}\par}
{\small\noindent\code{NR.transferBlock_refinement}\par}
{\small\noindent\code{NR.transferBlock_refinement_and_reconstruction}\par}
\smallskip\textbf{Checked scope.} Source and target projector families may live on different modules and refine independently; every coarse linear-map block is the double finite sum of its fine descendant blocks. The direct traceability wrapper also records reconstruction of the same map from all coarse blocks and from all fine blocks. The endomorphism formula is a specialization.
\par\smallskip\textbf{Boundary.} No sparse extraction algorithm, infinite block matrix, or norm convergence is proved here.
\medskip
\Needspace{7\baselineskip}
\subsection{Row 6: Residue projector algebra}
\label{row:6}
\textbf{Classification:} Exact.
\par\smallskip\textbf{Declaration targets}\par
{\small\noindent\code{R.residueProj_idempotent}\par}
{\small\noindent\code{R.residueProj_orthogonal}\par}
{\small\noindent\code{R.residueProj_complete}\par}
\smallskip\textbf{Checked scope.} For every positive modulus, coefficient residue selectors on \code{Nat -> A} are idempotent, mutually annihilating, and complete over an arbitrary semiring and module.
\par\smallskip\textbf{Boundary.} This is the division-free coefficient selector, not a character-average formula.
\medskip
\Needspace{7\baselineskip}
\subsection{Row 7: Root-of-unity realization}
\label{row:7}
\textbf{Classification:} Stronger / certificate-explicit.
\par\smallskip\textbf{Declaration targets}\par
{\small\noindent\code{R.ringFourierResidueKernel}\par}
{\small\noindent\code{R.ringFourierKernel_eq_residueKernel}\par}
{\small\noindent\code{R.FiniteRingFourierOrthogonalityCertificate}\par}
{\small\noindent\code{R.FiniteRingFourierOrthogonalityCertificate.toRingCertificate}\par}
{\small\noindent\code{R.ringFourierAverage}\par}
{\small\noindent\code{R.ringFourierAverage_eq_residueProj_of_finite_certificate}\par}
{\small\noindent\code{R.primitiveRoot_fourierOrthogonalityCertificate}\par}
{\small\noindent\code{R.fourierAverage_eq_residueProj_of_primitiveRoot}\par}
\smallskip\textbf{Checked scope.} Over any commutative semiring, \code{omega^N = 1} reduces the arbitrary coefficient address to a finite residue address. Consequently an explicit finite \code{N x N} orthogonality table identifies the division-free Fourier average with the residue selector. Over a field, a positive-order primitive root automatically supplies the normalized inverse-root certificate.
\par\smallskip\textbf{Boundary.} The general semiring branch still requires the finite orthogonality table; root-of-unity order alone is not enough in rings with zero divisors. The primitive-root construction is field-only.
\medskip
\Needspace{7\baselineskip}
\subsection{Row 8: Polyphase refinement}
\label{row:8}
\textbf{Classification:} Exact.
\par\smallskip\textbf{Declaration targets}\par
{\small\noindent\code{R.polyphase_refinement}\par}
\smallskip\textbf{Checked scope.} For positive \code{N} and \code{k}, the modulus-\code{N} selector is the finite sum of exactly those modulus-\code{N*k} selectors that merge to it, over any semiring module.
\par\smallskip\textbf{Boundary.} The theorem concerns exact coefficient masks; it gives no approximation or discarded-packet estimate.
\medskip
\Needspace{7\baselineskip}
\subsection{Row 9: Dyadic-odd packet factorization}
\label{row:9}
\textbf{Classification:} Exact.
\par\smallskip\textbf{Declaration targets}\par
{\small\noindent\code{R.crtEquiv}\par}
{\small\noindent\code{R.crtComponents_bijective}\par}
{\small\noindent\code{R.residueProj_crt}\par}
{\small\noindent\code{R.residueProj_crt_comm}\par}
{\small\noindent\code{R.dyadicOddEquiv}\par}
{\small\noindent\code{R.dyadicOddComponents_bijective}\par}
{\small\noindent\code{R.residueProj_dyadicOdd}\par}
\smallskip\textbf{Checked scope.} Positive coprime moduli give a bijection of residue labels and the composite selector in both composition orders. For positive odd \code{M}, this specializes to \code{2^a * M}.
\par\smallskip\textbf{Boundary.} This is coordinate CRT factorization. It does not identify the denominator or lifting behavior of spectral projectors.
\medskip
\Needspace{7\baselineskip}
\subsection{Row 10: Certified weighted inverse}
\label{row:10}
\textbf{Classification:} Stronger / certificate-explicit.
\par\smallskip\textbf{Declaration targets}\par
{\small\noindent\code{WT.TwoSidedInverse}\par}
{\small\noindent\code{WT.finiteGeometricInverse_spec_of_pow_eq_zero}\par}
{\small\noindent\code{WT.diagonalTransferInverse_spec}\par}
{\small\noindent\code{WT.analyticGeometricInverse_spec}\par}
{\small\noindent\code{WT.analyticDiagonalTransferInverse_spec}\par}
\smallskip\textbf{Checked scope.} In a possibly noncommutative ring, an explicit two-sided inverse for \code{D} and nilpotence of the normalized perturbation \code{Dinv*H} give a finite two-sided inverse of \code{D+H}. In a complete normed ring, \code{norm (Dinv*H) < 1} gives the separate infinite Neumann inverse.
\par\smallskip\textbf{Boundary.} The manuscript perturbation is \code{H}; the Lean declaration's parameter name \code{K} has no fraction-field meaning. There is no unrestricted convergence theorem for an abstract complete filtered ring. The analytic candidate is noncomputable.
\medskip
\Needspace{7\baselineskip}
\subsection{Row 11: Finite precision and weighted pruning}
\label{row:11}
\textbf{Classification:} Stronger / certificate-explicit.
\par\smallskip\textbf{Declaration targets}\par
{\small\noindent\code{WT.MultiplicativeFiltration}\par}
{\small\noindent\code{WT.PrecisionCertificate}\par}
{\small\noindent\code{WT.retainedFixedWordSum}\par}
{\small\noindent\code{WT.prunedWordInverse}\par}
{\small\noindent\code{WT.map_finiteGeometricInverse_eq_prunedWordInverse}\par}
{\small\noindent\code{WT.minWeight_mul_length_le_wordWeight}\par}
{\small\noindent\code{WT.map_sum_pow_eq_zero_at_scaled_cutoff}\par}
{\small\noindent\code{WT.prunedWordInverse_spec_of_minWeight}\par}
{\small\noindent\code{WT.finitePrecisionPrunedInverse_of_weightedEdges}\par}
{\small\noindent\code{WT.finitePrecisionPrunedDiagonalTransferInverse_of_minWeight}\par}
{\small\noindent\code{WT.finitePrecisionPrunedDiagonalTransferInverse_of_weightedEdges}\par}
\smallskip\textbf{Checked scope.} The mapped finite geometric inverse equals an explicit noncommutative word sum retaining exactly the words of cumulative weight below \code{L}. If every edge weight is at least \code{m} and \code{L <= m*N}, the \code{N}th mapped power vanishes and the pruned word sum is a two-sided inverse; the corresponding diagonal-plus-transfer formula follows from supplied \code{Dinv} and \code{Dinv*H = sum edge}.
\par\smallskip\textbf{Boundary.} The filtration, precision quotient map, edge-membership data, diagonal inverse, and edge expansion remain inputs. No DVR quotient, quotient map, basis-level modular matrix routine, or extracted pruning implementation is constructed.
\medskip
\Needspace{7\baselineskip}
\subsection{Row 12: Finite-graph transfer-zeta rationality}
\label{row:12}
\textbf{Classification:} Automatic.
\par\smallskip\textbf{Declaration targets}\par
{\small\noindent\code{TZ.coefficientAlgorithm_correct}\par}
{\small\noindent\code{TZ.coefficientAlgorithm_firstEdge}\par}
{\small\noindent\code{TZ.weightedAdjacency}\par}
{\small\noindent\code{TZ.transferMatrix}\par}
{\small\noindent\code{TZ.transferDenominator}\par}
{\small\noindent\code{TZ.transferNumerator}\par}
{\small\noindent\code{TZ.transferDenominator_coeff_zero}\par}
{\small\noindent\code{TZ.pathZeta_eq_matrixPathZeta}\par}
{\small\noindent\code{TZ.transferDenominator_mul_pathZeta}\par}
{\small\noindent\code{TZ.finite_graph_pathZeta_rationality}\par}
{\small\noindent\code{TZ.finite_graph_transferZeta_trace}\par}
{\small\noindent\code{TZC.certificate_explicit_rationality}\par}
\smallskip\textbf{Checked scope.} For every finite positive-weight directed multigraph and source/target pair, finite walk enumeration satisfies the first-edge recurrence. Lean constructs \code{Q = det(I-A(X))} and \code{P} as the corresponding adjugate entry, proves \code{Q(0)=1}, identifies the independently enumerated \code{pathZeta} with the matrix inverse entry, and proves \code{Q*Z=P} without a rationality certificate. A direct wrapper packages the manuscript's enumeration, recurrence, normalization, nonvanishing, and rational identity in one declaration. The older coefficient-certificate route remains available.
\par\smallskip\textbf{Boundary.} Walks may repeat vertices and edges; no simple-path restriction is imposed. Parallel edges are distinct edge values. Determinant and adjugate are noncomputable Lean definitions here; no denominator cancellation, minimality, extracted symbolic implementation, or analytic convergence claim is made.
\medskip
\Needspace{7\baselineskip}
\subsection{Row 13: Linear recurrence}
\label{row:13}
\textbf{Classification:} Automatic.
\par\smallskip\textbf{Declaration targets}\par
{\small\noindent\code{TZ.finite_graph_constant_coefficient_recurrence}\par}
{\small\noindent\code{TZC.constant_coefficient_recurrence}\par}
\smallskip\textbf{Checked scope.} The automatically constructed determinant denominator gives an antidiagonal constant-coefficient recurrence for every \code{n} beyond the degree of the canonical adjugate numerator. A separately supplied \code{P,Q} certificate yields the same form through the legacy route.
\par\smallskip\textbf{Boundary.} The canonical recurrence is not claimed minimal, and no recurrence is guessed or certified from finitely many sampled coefficients alone.
\medskip
\Needspace{7\baselineskip}
\subsection{Row 14: Evaluated Bezout spectral projector}
\label{row:14}
\textbf{Classification:} Stronger / certificate-explicit.
\par\smallskip\textbf{Declaration targets}\par
{\small\noindent\code{RS.EvaluatedBezoutCertificate}\par}
{\small\noindent\code{RS.separatorNumerator_eq_scalar_sub_of_bezout}\par}
{\small\noindent\code{RS.separatorNumerator_mul_self_of_bezout}\par}
{\small\noindent\code{RS.separator_idempotent_of_bezout}\par}
{\small\noindent\code{RS.separatorNumerator_apply_of_mem_ker_F_of_bezout}\par}
{\small\noindent\code{RS.separator_apply_of_mem_ker_F_of_bezout}\par}
{\small\noindent\code{RS.separator_apply_of_mem_ker_G}\par}
{\small\noindent\code{RS.scalar_mul_separator_of_ne_zero}\par}
{\small\noindent\code{RS.scalar_mul_separator_preservesLattice_of_ne_zero}\par}
\smallskip\textbf{Checked scope.} Over a field, the evaluated equality \code{U*F + V*G = r*1}, \code{F*G=0}, and only the two commutations actually used prove the denominator-free identity \code{(VG)^2=r(VG)}; adding \code{r!=0} gives idempotence of \code{r^-1*VG}. The kernel actions need only the displayed evaluated Bezout or kernel hypotheses, and lattice preservation of the numerator follows when \code{V,G} preserve the lattice.
\par\smallskip\textbf{Boundary.} The evaluated Bezout equality and annihilation are supplied. No polynomial Bezout identity, resultant, characteristic factorization, or Cayley-Hamilton instance is constructed; nilpotence alone is not used to infer exact kernel action.
\medskip
\Needspace{7\baselineskip}
\subsection{Row 15: Separation weight}
\label{row:15}
\textbf{Classification:} Stronger / certificate-explicit.
\par\smallskip\textbf{Declaration targets}\par
{\small\noindent\code{RS.ScalarPreservesLattice}\par}
{\small\noindent\code{RS.IsLatticeUnit}\par}
{\small\noindent\code{RS.IsSeparationWeight}\par}
{\small\noindent\code{RS.IsSeparationWeight.mono}\par}
{\small\noindent\code{RS.separationWeight}\par}
{\small\noindent\code{RS.separationWeight_clears}\par}
{\small\noindent\code{RS.separationWeight_le_of_clears}\par}
{\small\noindent\code{RS.pow_smul_separator_eq_unit_inv_smul_separatorNumerator}\par}
{\small\noindent\code{RS.separator_isSeparationWeight_of_unit_mul_pow_eq_of_inv_preserves}\par}
{\small\noindent\code{RS.separator_hasFiniteSeparationWeight_of_eq_unit_mul_pow}\par}
{\small\noindent\code{RS.separator_separationWeight_le_of_eq_unit_mul_pow}\par}
\smallskip\textbf{Checked scope.} \code{IsSeparationWeight pi E lattice s} means that \code{(pi^s)E} preserves the chosen lattice; weights are upward closed when scalar multiplication by \code{pi} preserves it. If \code{r=u*pi^s}, then \code{pi^s E=u^-1 VG}; preservation by the inverse unit and numerator proves clearing, while a bundled lattice-unit certificate gives existence and the least-weight upper bound \code{separationWeight<=s}.
\par\smallskip\textbf{Boundary.} The unit action on the chosen lattice is explicit data. The least weight is noncomputable via \code{Nat.find}; no valuation, polynomial resultant construction, search procedure, or claim that the displayed exponent is minimal is supplied.
\medskip
\Needspace{7\baselineskip}
\subsection{Row 16: Projector precision delay}
\label{row:16}
\textbf{Classification:} Exact from certificate.
\par\smallskip\textbf{Declaration targets}\par
{\small\noindent\code{RS.CongruentModLattice}\par}
{\small\noindent\code{RS.projector_precision_delay}\par}
\smallskip\textbf{Checked scope.} For an arbitrary field scalar \code{pi}, additive lattice, and clearing-weight certificate, an explicit congruence witness at depth \code{m+s} maps to one at depth \code{m}.
\par\smallskip\textbf{Boundary.} Congruence is witness-based scalar-lattice membership; no topology, completion, digit representation, or valuation API is required or concluded.
\medskip
\Needspace{7\baselineskip}
\subsection{Row 17: Odd-adic versus dyadic sign behavior}
\label{row:17}
\textbf{Classification:} Stronger / scoped realization.
\par\smallskip\textbf{Declaration targets}\par
{\small\noindent\code{SD.signInvolution_mul_self}\par}
{\small\noindent\code{SD.positiveSignProjector_idempotent}\par}
{\small\noindent\code{SD.negativeSignProjector_idempotent}\par}
{\small\noindent\code{SD.signProjectors_complete}\par}
{\small\noindent\code{SD.HasExactPLocalDenominatorWeight}\par}
{\small\noindent\code{SD.positiveSignProjector_exact_pLocal_weight_zero_of_prime_ne_two}\par}
{\small\noindent\code{SD.negativeSignProjector_exact_pLocal_weight_zero_of_prime_ne_two}\par}
{\small\noindent\code{SD.signFactors_resultant}\par}
{\small\noindent\code{SD.signProjectors_resultant_dyadic_weight_endpoint}\par}
{\small\noindent\code{SD.padicValNat_two_at_odd_prime}\par}
{\small\noindent\code{SD.oddAdicVersusDyadicSignBehavior_trace}\par}
\smallskip\textbf{Checked scope.} The displayed rational \code{2 x 2} involution and sign projectors are checked entrywise. Each projector has exact local denominator weight zero at every prime other than \code{2} and exact \code{2}-local weight one. Lean also computes \code{resultant(X-1,X+1)=2} and identifies its \code{2}-adic valuation with the exact local weight of both projectors. The direct trace wrapper combines both local regimes and the resultant statements.
\par\smallskip\textbf{Boundary.} This is an explicit rational-entry, natural-valuation model. Reduction modulo \code{2} applies to the characteristic factors, not to the rational projector matrices containing \code{1/2}. Lean does not construct \code{Q_p}, its valuation ring, or a general equivalence between rational \code{p}-integrality and preservation of an arbitrary completed lattice.
\medskip
\Needspace{7\baselineskip}
\subsection{Row 18: Concrete ZMod Smith-weighted congruence solve}
\label{row:18}
\textbf{Classification:} Automatic at concrete scope / certificate-explicit extension.
\par\smallskip\textbf{Declaration targets}\par
{\small\noindent\code{SD.smithZModMulHom}\par}
{\small\noindent\code{SD.SmithZModAdmissible}\par}
{\small\noindent\code{SD.smithZModFiber_nonempty_iff}\par}
{\small\noindent\code{SD.smithZModFiber_card_of_admissible}\par}
{\small\noindent\code{SD.smithZModCoordinateCertificate}\par}
{\small\noindent\code{SD.smithZMod_global_solvable_iff}\par}
{\small\noindent\code{SD.smithZMod_global_conditional_count}\par}
{\small\noindent\code{SD.SmithSolvabilityCertificate}\par}
{\small\noindent\code{SD.SmithCardinalityCertificate}\par}
{\small\noindent\code{SD.smith_conditional_product_count}\par}
\smallskip\textbf{Checked scope.} For \code{p>1}, Lean proves directly over \code{ZMod(p^depth)} that \code{p^weight*x=target} is solvable exactly when \code{p^min(weight,depth)} divides \code{target.val}, and every admissible fibre has cardinality \code{p^min(weight,depth)}. It then proves the coordinatewise criterion and the exact conditional count \code{p^(sum min(weight,depth))} for finite diagonal systems, including cardinality zero when a coordinate is inadmissible. The split certificate API gives the analogous product statement for abstract finite fibres, with the residue cardinality supplied as \code{q}.
\par\smallskip\textbf{Boundary.} The counted manuscript theorem is the concrete \code{ZMod(p^L)} result with \code{q=p}. The arbitrary-DVR, residue-size-\code{q} extension is classical/certificate-explicit, not derived by Lean. Diagonal weights and targets are inputs; Lean does not construct Smith normal form. A literal zero diagonal factor is represented at precision \code{L} by a weight at least \code{L}, or handled by an explicit abstract coordinate certificate.
\medskip
\Needspace{7\baselineskip}
\subsection{Row 19: Power-defect identity}
\label{row:19}
\textbf{Classification:} Stronger.
\par\smallskip\textbf{Declaration targets}\par
{\small\noindent\code{TLD.defect}\par}
{\small\noindent\code{TLD.power_defect_identity_all}\par}
{\small\noindent\code{TLD.power_defect_identity}\par}
{\small\noindent\code{TLD.powers_compatible_of_defect_eq_zero}\par}
{\small\noindent\code{TLD.polynomialAction}\par}
{\small\noindent\code{TLD.polynomialAction_compatible_of_defect_eq_zero}\par}
{\small\noindent\code{TLD.continuousPolynomialAction_compatible_of_defect_eq_zero}\par}
{\small\noindent\code{LD.power_defect_identity_all}\par}
\smallskip\textbf{Checked scope.} The primary theorem is typed for distinct fine and coarse semiring modules and proves the ordered telescoping identity for every natural \code{n}, including zero. Zero one-step defect gives exact compatibility with every power and every finite polynomial functional calculus; the continuous version is also proved.
\par\smallskip\textbf{Boundary.} This is linear algebra and finite polynomial calculus; no nonlinear higher-jet, operator-tree defect expansion, or analytic functional calculus is asserted.
\medskip
\Needspace{7\baselineskip}
\subsection{Row 20: Power-defect norm bound}
\label{row:20}
\textbf{Classification:} Stronger.
\par\smallskip\textbf{Declaration targets}\par
{\small\noindent\code{TLD.continuousDefect}\par}
{\small\noindent\code{TLD.continuous_power_defect_identity_all}\par}
{\small\noindent\code{TLD.continuous_end_pow_norm_le}\par}
{\small\noindent\code{TLD.continuous_norm_power_defect_le}\par}
{\small\noindent\code{TLD.continuous_norm_power_defect_le_geometric_majorant}\par}
{\small\noindent\code{LD.norm_power_defect_le_raw}\par}
{\small\noindent\code{LD.norm_power_defect_le}\par}
\smallskip\textbf{Checked scope.} Over a nontrivially normed field, continuous maps between distinct normed fine and coarse spaces satisfy the exact finite-sum operator-norm bound for every natural \code{n}, including zero. If both dynamics have norm at most \code{rho}, Lean derives the closed majorant \code{n*rho^(n-1)*norm(defect)}. No completeness assumption on the spaces is needed.
\par\smallskip\textbf{Boundary.} These are operator-norm estimates, not spectral-radius estimates, equality statements, or nonlinear defect bounds. The majorant takes \code{rho} and both norm inequalities as hypotheses.
\medskip
\Needspace{7\baselineskip}
\subsection{Row 21: Resolvent-defect identity}
\label{row:21}
\textbf{Classification:} Stronger / certificate-explicit.
\par\smallskip\textbf{Declaration targets}\par
{\small\noindent\code{TLD.TwoSidedInverse}\par}
{\small\noindent\code{TLD.scalarShift_resolvent_identity}\par}
{\small\noindent\code{TLD.resolvents_compatible_of_defect_eq_zero}\par}
{\small\noindent\code{TLD.continuousScalarShift}\par}
{\small\noindent\code{TLD.continuous_scalarShift_resolvent_identity}\par}
{\small\noindent\code{TLD.continuous_norm_resolvent_defect_le}\par}
{\small\noindent\code{TLD.continuous_resolvents_compatible_of_defect_eq_zero}\par}
{\small\noindent\code{LD.resolvent_defect_identity}\par}
\smallskip\textbf{Checked scope.} Algebraically, supplied two-sided inverse witnesses give the typed resolvent identity on distinct modules. For continuous maps, the proof needs only \code{(zI-Tf)Rf=I} and \code{Rc(zI-Tc)=I} and yields \code{norm(C Rf-Rc C) <= norm(Rc)*norm(defect)*norm(Rf)}; zero defect gives exact intertwining.
\par\smallskip\textbf{Boundary.} All inverse equations are supplied witnesses. Lean does not construct resolvents, prove spectral membership, or infer inverse existence from analytic hypotheses.
\medskip
\Needspace{7\baselineskip}
\subsection{Row 22: Smith phase inversion identity}
\label{row:22}
\textbf{Classification:} Stronger / scoped realization.
\par\smallskip\textbf{Declaration targets}\par
{\small\noindent\code{SD.smithPhase}\par}
{\small\noindent\code{SD.smithPhase_neg_one_eq_inv}\par}
\smallskip\textbf{Checked scope.} Over any field in which \code{2} is nonzero, the pulled-back phase formula at \code{omega = -1} equals \code{z^-1} for \code{z != 0}; this includes \code{Q}, \code{R}, and \code{C} and does not need the manuscript's additional \code{z != -1} condition.
\par\smallskip\textbf{Boundary.} Only the algebraic formula is checked. No Smith disk, Mobius-domain, radius-preservation, or physical impedance theorem is formalized.
\medskip
\Needspace{7\baselineskip}
\subsection{Row 23: Residue analysis and synthesis algorithm}
\label{row:23}
\textbf{Classification:} Specification.
\par\smallskip\textbf{Declaration targets}\par
{\small\noindent\code{PC.rangeSynthesis_rangeAnalysis}\par}
{\small\noindent\code{PC.rangeAnalysis_rangeSynthesis}\par}
{\small\noindent\code{PC.exactPacketCoordinates}\par}
{\small\noindent\code{NR.coordinateRefinementEquiv}\par}
\smallskip\textbf{Checked scope.} For finite complete projector families, range-restricted analysis and finite-sum synthesis are inverse linear maps; cross-level analysis/synthesis gives an exact coordinate equivalence.
\par\smallskip\textbf{Boundary.} This is the mathematical contract behind the array routine, not extracted interleaving code or a proof about Python storage shapes.
\medskip
\Needspace{7\baselineskip}
\subsection{Row 24: Weighted finite-precision inverse algorithm}
\label{row:24}
\textbf{Classification:} Specification / certificate-explicit.
\par\smallskip\textbf{Declaration targets}\par
{\small\noindent\code{WT.finiteGeometricInverse}\par}
{\small\noindent\code{WT.prunedWordInverse}\par}
{\small\noindent\code{WT.map_finiteGeometricInverse_eq_prunedWordInverse}\par}
{\small\noindent\code{WT.le_mul_ceilDiv}\par}
{\small\noindent\code{WT.prunedWordInverse_spec_of_minWeight}\par}
{\small\noindent\code{WT.prunedDiagonalTransferInverse}\par}
{\small\noindent\code{WT.finitePrecisionPrunedDiagonalTransferInverse_of_minWeight}\par}
{\small\noindent\code{WT.finitePrecisionPrunedDiagonalTransferInverse_of_weightedEdges}\par}
\smallskip\textbf{Checked scope.} The bounded geometric loop maps exactly to an explicit signed sum of retained noncommutative words. For \code{m>0}, Lean proves \code{L <= m * ceil(L/m)}, so the pseudocode choice \code{N=ceil(L/m)} meets the cutoff premise. Filtration data and \code{L<=m*N} prove that the pruned expression is a two-sided inverse; supplied \code{Dinv} and a normalized-transfer edge expansion lift it to the explicit pruned diagonal-plus-transfer inverse after applying \code{q}.
\par\smallskip\textbf{Boundary.} This is a mathematical specification. The manuscript calls the perturbation \code{H}; the Lean parameter name \code{K} has no fraction-field meaning. The result does not provide an extracted modular matrix inverse, basis representation, zero test, complexity bound, or proof of equivalence to the companion Python implementation.
\medskip
\Needspace{7\baselineskip}
\subsection{Row 25: Transfer-zeta walk construction algorithm}
\label{row:25}
\textbf{Classification:} Specification / automatic.
\par\smallskip\textbf{Declaration targets}\par
{\small\noindent\code{TZ.walkTuplesAtWeight}\par}
{\small\noindent\code{TZ.mem_walkTuplesAtWeight_iff}\par}
{\small\noindent\code{TZ.tuple_captured_iff}\par}
{\small\noindent\code{TZ.coefficientAlgorithm}\par}
{\small\noindent\code{TZ.coefficientAlgorithm_correct}\par}
{\small\noindent\code{TZ.exactWalkFirstEdgeEquiv}\par}
{\small\noindent\code{TZ.coefficientAlgorithm_firstEdge}\par}
{\small\noindent\code{TZ.pathZetaAlgorithm_trace}\par}
{\small\noindent\code{TZ.weightedAdjacency}\par}
{\small\noindent\code{TZ.weightedAdjacency_coeff}\par}
{\small\noindent\code{TZ.enumeratedPathZetaMatrix_eq_matrixPathZetaMatrix}\par}
{\small\noindent\code{TZ.pathZeta_eq_inverse_mul_transferNumerator}\par}
{\small\noindent\code{TZ.finite_graph_pathZeta_rationality}\par}
{\small\noindent\code{TZ.finite_graph_constant_coefficient_recurrence}\par}
\smallskip\textbf{Checked scope.} The coefficient loop is a finite \code{Finset} enumeration with exact capture and cardinality theorems. A direct trace wrapper pairs its semantic walk count with the first-edge recurrence. Finite matrix determinant/adjugate data are then constructed and proved equal to the enumerated series, yielding the normalized rational identity and eventual recurrence without a supplied rationality certificate.
\par\smallskip\textbf{Boundary.} The objects counted are directed walks, with repetitions allowed. The Lean polynomial matrix definitions are noncomputable, and no extracted determinant/adjugate program, cancellation or minimal-denominator routine, performance bound, or analytic convergence theorem is claimed.
\medskip

\clearpage
\appendix
\section{Final traditional-paper navigation}\label{app:crosswalk}
The printed page numbers below refer to the final traditional companion, whose title page is unnumbered. They are not PDF viewer page indices. A shared label is a navigation aid; it does not extend the formal scope to additional claims in the surrounding section. In particular, the Smith-phase display has no equation number.
\begin{longtable}{p{.07\textwidth}p{.50\textwidth}p{.32\textwidth}}
\toprule Row & Canonical label & Final locator \\ \midrule\endhead
1 & \code{def:refinement-map} & Statement 3.2, p.~7 \\ \addlinespace
2 & \code{thm:nested-refinement} & Statement 3.3, p.~7 \\ \addlinespace
3 & \code{def:coordinate-refinement} & Statement 3.4, p.~8 \\ \addlinespace
4 & \code{prop:perfect-reconstruction} & Statement 3.5, p.~8 \\ \addlinespace
5 & \code{thm:transfer-block-refinement} & Statement 3.8, p.~9 \\ \addlinespace
6 & \code{prop:residue-algebra} & Statement 4.2, p.~14 \\ \addlinespace
7 & \code{thm:root-of-unity} & Statement 4.3, p.~14 \\ \addlinespace
8 & \code{thm:polyphase-refinement} & Statement 4.6, p.~15 \\ \addlinespace
9 & \code{thm:dyadic-odd} & Statement 4.9, p.~16 \\ \addlinespace
10 & \code{thm:weighted-inverse} & Statement 5.10, p.~21 \\ \addlinespace
11 & \code{cor:finite-precision} & Statement 5.11, p.~21 \\ \addlinespace
12 & \code{thm:transfer-zeta} & Statement 6.6, p.~26 \\ \addlinespace
13 & \code{cor:zeta-recurrence} & Statement 6.7, p.~26 \\ \addlinespace
14 & \code{thm:bezout-projector} & Statement 7.5, p.~30 \\ \addlinespace
15 & \code{def:separation-weight} & Statement 7.6, p.~30 \\ \addlinespace
16 & \code{thm:precision-delay} & Statement 7.7, p.~31 \\ \addlinespace
17 & \code{prop:dyadic-sign} & Statement 9.3, p.~37 \\ \addlinespace
18 & \code{thm:smith-solve} & Statement 10.6, p.~41 \\ \addlinespace
19 & \code{thm:power-defect} & Statement 11.1, p.~42 \\ \addlinespace
20 & \code{cor:defect-norm} & Statement 11.3, p.~43 \\ \addlinespace
21 & \code{thm:resolvent-defect} & Statement 11.5, p.~43 \\ \addlinespace
22 & \code{eq:smith-inversion} & Section 12.3, Smith-phase identity, p.~49 \\ \addlinespace
23 & \code{alg:residue} & Statement B.1, p.~53 \\ \addlinespace
24 & \code{alg:weighted-inverse} & Statement B.2, p.~53 \\ \addlinespace
25 & \code{alg:path-zeta} & Statement B.3, p.~53 \\ \addlinespace
\bottomrule\end{longtable}

\section{Source imports and pinned dependencies}\label{app:imports}
This table records direct local imports. It is a file dependency graph, not a theorem-level logical dependency analysis. The JSON companion also retains every direct external import. Some modules import the broad \code{Mathlib} module, while others list focused imports.
\begin{longtable}{p{.31\textwidth}p{.59\textwidth}}
\toprule Module & Direct local imports \\ \midrule\endhead
\code{ExactCoordinates} & \code{ProjectorCore} \\ \addlinespace
\code{FibreRefinement} & \code{ProjectorCore} \\ \addlinespace
\code{PaperII.KernelAudit} & \code{PaperII.NestedRefinement}, \code{PaperII.ResidueFourierCRT}, \code{PaperII.WeightedTransfer}, \code{PaperII.TransferZeta}, \code{PaperII.ResultantSeparation}, \code{PaperII.SmithDyadic}, \code{PaperII.LinearDefects}, \code{PaperII.TypedLinearDefects} \\ \addlinespace
\code{PaperII.LinearDefects} & External imports only. \\ \addlinespace
\code{PaperII.NestedRefinement} & \code{ExactCoordinates} \\ \addlinespace
\code{PaperII.ResidueFourierCRT} & \code{ResidueProjectors} \\ \addlinespace
\code{PaperII.ResultantSeparation} & External imports only. \\ \addlinespace
\code{PaperII.SmithDyadic} & External imports only. \\ \addlinespace
\code{PaperII.TransferZeta} & External imports only. \\ \addlinespace
\code{PaperII.TypedLinearDefects} & External imports only. \\ \addlinespace
\code{PaperII.WeightedTransfer} & External imports only. \\ \addlinespace
\code{PaperII} & \code{PaperII.NestedRefinement}, \code{PaperII.ResidueFourierCRT}, \code{PaperII.WeightedTransfer}, \code{PaperII.TransferZeta}, \code{PaperII.ResultantSeparation}, \code{PaperII.SmithDyadic}, \code{PaperII.LinearDefects}, \code{PaperII.TypedLinearDefects} \\ \addlinespace
\code{ProjectorCore} & External imports only. \\ \addlinespace
\code{ResidueProjectors} & \code{FibreRefinement} \\ \addlinespace
\code{KomorebiVerification} & \code{PaperII} \\ \addlinespace
\bottomrule\end{longtable}

The exact toolchain string is \code{leanprover/lean4:v4.34.0-rc2}. The Lake manifest pins the following source revisions. Its descriptive \code{inputRev} fields such as \code{master} and \code{main} do not replace these exact revisions.
\begin{longtable}{p{.21\textwidth}p{.69\textwidth}}
\toprule Package & Exact revision \\ \midrule\endhead
mathlib & \code{aeb72f1e9981b6d35b426b16dca7759be3e7ba06} \\ \addlinespace
plausible & \code{d9598f07b1bc701f1e3aae163d2681c1fd978793} \\ \addlinespace
LeanSearchClient & \code{ba67e212be1197b84c1f1f6299488a10a3002713} \\ \addlinespace
importGraph & \code{d8823026ac7ef130c253089d95685f9877b95323} \\ \addlinespace
proofwidgets & \code{a8acbfd87375ff4abe14ce09db5b7664d383bc7f} \\ \addlinespace
aesop & \code{18889deb9e83ea7420ef51c160d6f88552e744e3} \\ \addlinespace
Qq & \code{507746ab8f4b643ccdacb2ec4cdb5853fa9f8ab3} \\ \addlinespace
batteries & \code{4cac2177c37f5530c4da76aa8e4307f3fc9e4dcb} \\ \addlinespace
Cli & \code{ab3a82db9fea14cf0fd7f5a2de650f4b534640af} \\ \addlinespace
\bottomrule\end{longtable}

\section{Source-excerpt and reproduction inventory}\label{app:inventory}
Every printed Lean listing comes from an exact UTF-8 excerpt in \code{excerpts/*.lean.txt}. The adjacent \code{*.render.txt} file escapes TeX control characters and substitutes a math-font glyph where the code font lacks one. It preserves token order. Exact source text, not the display fallback, is the object to copy into an editor. Excerpts retain their surrounding source context only when that context is printed; they are not standalone Lean files.
\code{generated/excerpts.json} records each source path, inclusive source line range, full-source SHA-256, exact-excerpt SHA-256, and display-file SHA-256. \code{generated/theorem_map.json} contains all 25 rows with fully qualified targets, classifications, scopes, boundaries, and lexical source locations. \code{generated/mapped_source_index.json} provides an alphabetical target index. \code{generated/static_identity_audit.json} records all manifest comparisons and the optional baseline comparison performed for this edition.

\code{build_guide_data.py} uses only the Python standard library. It generates the map, import table, pins, excerpts, and current static audit without invoking Lean. \code{build_guide.sh} calls that generator and runs XeLaTeX four times. The final PDF requires the P052, DejaVu Sans, DejaVu Sans Mono, and Latin Modern Math fonts used by this edition, together with the LaTeX packages listed in its preamble. The independent PDF check script uses PyMuPDF and Pillow for page rendering and contact sheets.

The source project remains unchanged. Editing this guide cannot by itself repair or strengthen a Lean theorem. A mathematical change to the formal component requires a separate source change, a new compile, an updated map when appropriate, and new execution evidence.
\end{document}
